\section*{Capítulo \ref{ch:g8:sound-pitch-loudness} --- El sonido: velocidad, tono e intensidad sonora}

\begin{solution}{exo:g8:sound-pitch-loudness:1}
\omterm{def:g8:sound-pitch-loudness:frequency}{Frecuencia}: el número de vaivenes completos por segundo, en
\omterm{def:g8:sound-pitch-loudness:frequency}{hercios} (\unit{Hz}). El la de concierto: \qty{440}{Hz}; la octava
de encima: \qty{880}{Hz} --- el doble.
\end{solution}

\begin{solution}{exo:g8:sound-pitch-loudness:2}
$4$ vaivenes en \qty{0.02}{s}: \qty{200}{Hz}; $16$ vaivenes:
\qty{800}{Hz} --- el segundo suena más agudo (dos octavas, de
hecho).
\end{solution}

\begin{solution}{exo:g8:sound-pitch-loudness:3}
De unos \qty{20}{Hz} a \qty{20000}{Hz}. Por debajo: los
\omterm{prop:g8:sound-pitch-loudness:range}{infrasonidos} --- los retumbos de largo alcance de los elefantes.
Por encima: los \omterm{prop:g8:sound-pitch-loudness:range}{ultrasonidos} --- los chasquidos de caza del
murciélago.
\end{solution}

\begin{solution}{exo:g8:sound-pitch-loudness:4}
Más fuerte, la misma nota: los picos crecen en altura y el espaciado no
cambia. Más aguda, el mismo \omterm{def:g6:mass-and-volume:volume}{volumen}: los picos se apretujan y la altura
no cambia. Dos mandos independientes.
\end{solution}

\begin{solution}{exo:g8:sound-pitch-loudness:5}
Los tres suben la \omterm{def:g8:sound-pitch-loudness:frequency}{frecuencia}: las cuerdas más cortas, más tensas
y más delgadas tiemblan más veces por segundo --- $f$ más alta, tono
más agudo.
\end{solution}

\begin{solution}{exo:g8:sound-pitch-loudness:6}
Hojas \qty{10}{dB}; tráfico \qty{80}{dB}; primera fila \qty{110}{dB};
dolor cerca de \qty{120}{dB}. La banda del daño se abre alrededor de
los \qty{85}{dB}.
\end{solution}

\begin{solution}{exo:g8:sound-pitch-loudness:7}
El tono del silbato se sitúa justo por encima de los \qty{20000}{Hz}
--- fuera de la ventana del dueño y dentro de la del perro. Estropeado
para un par de oídos y penetrante para el otro: las ventanas cambian
según la especie.
\end{solution}

\begin{solution}{exo:g8:sound-pitch-loudness:8}
Que la \omterm{def:g7:motion-average-speed:formula}{velocidad} es una y la misma para todos los tonos y todos
los volúmenes: \omterm{def:g7:motion-average-speed:formula}{velocidades} distintas emborronarían los acordes
lejanos en arpegios y retrasarían lo flojo detrás de lo fuerte --- el
concierto del parque refuta las dos cosas cada noche.
\end{solution}

\begin{solution}{exo:g8:sound-pitch-loudness:9}
$50000 \times 0.01 = 500$ vaivenes por chasquido. La pintura fina
necesita pinceles finos: solo las vibraciones de vaivenes muy
cortos --- de $f$ muy alta --- pueden palpar detalles pequeños en sus
\omterm{ex:g4:how-sound-travels:echo}{ecos}; los retumbos toscos de baja \omterm{def:g8:sound-pitch-loudness:frequency}{frecuencia} pasan por
encima de una polilla sin enterarse.
\end{solution}

\begin{solution}{exo:g8:sound-pitch-loudness:10}
Las dos teclas están dentro de la ventana --- los \qty{27}{Hz} justo
por encima del suelo de \qty{20}{Hz} y los \qty{4200}{Hz} cómodamente
por debajo del techo. Pero la edad erosiona la ventana por
\emph{arriba}: el propio oído del afinador se embota primero en los
agudos, así que se llama al analizador para las octavas más altas antes
que para las más bajas.
\end{solution}

\begin{solution}{exo:g8:sound-pitch-loudness:11}
Un crecimiento por factores de diez escondido en pasos iguales --- cada
pocos peldaños multiplican la \omterm{def:g2:pushes-and-pulls:force}{fuerza} del \omterm{def:g3:sound-around-us:sound}{sonido}, no la suman. Así
que ``un pelín más alto'' en la escalera es \emph{mucho} más fuerte
para el oído: los presupuestos se derrumban de \omterm{ex:g2:measuring-time:clock}{horas} a minutos a
lo largo de lo que se lee como un giro pequeño del mando.
\end{solution}

\begin{solution}{exo:g8:sound-pitch-loudness:12}
Comprobación del piano: el analizador junto a la cuerda, y comparar su
lectura con $440$. \omterm{def:g6:mass-and-volume:volume}{Volumen} de la última fila: colocarse allí con la
lectura de \omterm{ex:g8:sound-pitch-loudness:decibels}{decibelios} durante la prueba de batería. Ventilación:
apuntar la pantalla de \omterm{prop:g8:colors-spectra:mixture}{espectro} del analizador al extremo grave y
buscar un pico por debajo de \qty{20}{Hz}. Lo menos fiable: el
veredicto sobre los \omterm{prop:g8:sound-pitch-loudness:range}{infrasonidos} --- los micrófonos de móvil están
construidos para voces y se quedan sordos cerca del suelo de la
ventana, así que una pantalla en silencio demuestra poco ahí abajo.
\end{solution}

\begin{solution}{pb:g8:sound-pitch-loudness:1}
\textbf{1.} El piano está bajo de afinación --- todos sus la tiemblan
ocho vaivenes por segundo de menos. El mando del tono: la llave de un
afinador, no un amplificador.
\textbf{2.} Los \qty{41}{Hz} se sitúan justo por encima del suelo de
\qty{20}{Hz} de la ventana: se oyen, y muy graves --- y con vaivenes
tan lentos, el propio tambor del pecho responde también: se nota tanto
como se oye.
\textbf{3.} El mismo tono (espaciado igual, \qty{440}{Hz}) y tres veces
la \omterm{prop:g8:sound-pitch-loudness:loudness}{amplitud}: el cantante canta la nota del diapasón, mucho más
fuerte.
\textbf{4.} Solo los oídos más jóvenes --- el techo de la ventana,
cerca de los \qty{20000}{Hz}, es de los niños; las ventanas de casi
todos los adultos ya han bajado, y ningún oyente certifica
\qty{22000}{Hz}.
\textbf{5.} La sesión de rock (\qty{100}{dB}) y el final
(\qty{110}{dB}) están las dos por encima de la línea de daño de los
\qty{85}{dB}: territorio de solo minutos; la sesión acústica pasa.
\textbf{6.} Unos \qty{80}{dB} detrás de los tapones --- otra vez por
debajo de la línea de daño: una exposición de toda la velada convertida
en soportable, que es la razón de que los sensatos los lleven siempre.
\textbf{7.} Las cáscaras que se extienden: cada apretujamiento se
reparte con la distancia por una cáscara más vasta --- el volumen (la
\omterm{prop:g8:sound-pitch-loudness:loudness}{amplitud}) se apaga. El tono viaja sin cambios: la distancia gira
un solo mando.
\textbf{8.} Tocar suave gira únicamente el mando de la
\omterm{prop:g8:sound-pitch-loudness:loudness}{amplitud} --- la \omterm{def:g8:sound-pitch-loudness:frequency}{frecuencia} de cada nota, aguda o
grave, se queda exactamente donde estaba. (Su ``apagado'' es una manía
del propio oído a volúmenes bajos, no la física de las cuerdas.)
\textbf{9.} $680 \div 340 = \qty{2}{s}$ del destello al estallido. Dos
golpes por segundo: \qty{2}{Hz} --- muy por debajo del suelo: se nota
en el pecho, no se oye nunca como un tono.
\textbf{10.} Los vaivenes más graves del bajo: cerca del suelo de la
ventana y por debajo, exactamente donde se queda sordo un micrófono
construido para voces --- el retumbo era real; el móvil no lo capturó
nunca.
\textbf{11.} Por ejemplo: ``El tono es la \omterm{def:g8:sound-pitch-loudness:frequency}{frecuencia} --- la cuenta
de vaivenes por segundo, y todo el oficio del afinador. El \omterm{def:g6:mass-and-volume:volume}{volumen} es
la \omterm{prop:g8:sound-pitch-loudness:loudness}{amplitud} --- el tamaño de los vaivenes, y todo el libro de
cuentas de la enfermera. Y ninguno de los dos mandos toca la
\omterm{def:g7:motion-average-speed:formula}{velocidad}: todas las notas de esta noche cruzaron el salón a los
mismos \qty{340}{m/s}''.
\textbf{12.} La proposición de la ventana: el ahuyentador afinado por
encima de \qty{20000}{Hz} --- fuera de todas las ventanas humanas y
dentro de la del ratón. Un solo temblor, dos audiencias, ninguna
contradicción.
\end{solution}
